\documentclass[UTF8,a4paper,12pt]{ctexart}

\usepackage{geometry}
\usepackage{booktabs}
\usepackage{longtable}
\usepackage{array}
\usepackage{graphicx}
\usepackage{float}
\usepackage{hyperref}
\usepackage{xcolor}
\usepackage{caption}

\geometry{left=2.5cm,right=2.5cm,top=2.8cm,bottom=2.8cm}
\graphicspath{{reports/figures/}}

\title{基于股票快照的价格移动方向预测：EDA 分析报告}
\author{}
\date{}

\begin{document}

\maketitle

\section{EDA 目标}

本次探索性数据分析（Exploratory Data Analysis, EDA）的目标是对股票快照订单簿数据进行系统检查，为后续的数据清洗、特征工程和建模提供依据。重点关注以下问题：

\begin{enumerate}
    \item 数据文件是否完整，是否存在读取失败、空文件或缺失交易组合；
    \item 字段结构是否符合项目说明，关键变量是否可用；
    \item 数据中是否存在缺失值、重复值、无穷值等基础质量问题；
    \item 订单簿价格和挂单量是否存在异常；
    \item \texttt{n\_midprice} 是否与买一价、卖一价规则基本一致；
    \item \texttt{label\_5}、\texttt{label\_10}、\texttt{label\_20}、\texttt{label\_40}、\texttt{label\_60} 是否对应未来价格移动方向；
    \item 标签类别是否均衡，后续建模是否需要类别权重或采样策略。
\end{enumerate}

\section{项目任务与数据背景}

本项目的任务是基于股票订单簿快照数据，预测未来若干 tick 后中间价的移动方向。预测窗口包括未来 5、10、20、40、60 个 tick。目标变量为三分类标签：0 表示下跌，1 表示不变，2 表示上涨。

实际数据中的标签字段为：

\begin{center}
\texttt{label\_5, label\_10, label\_20, label\_40, label\_60}
\end{center}

文件命名格式为：

\begin{center}
\texttt{snapshot\_sym<xx>\_date<yy>\_am.csv} 或
\texttt{snapshot\_sym<xx>\_date<yy>\_pm.csv}
\end{center}

每个文件表示某一只股票、某一个交易日、某一个交易时段的行情快照数据。

\section{数据概览}

本次 EDA 共读取行情文件 1521 个，全部读取成功。合并后数据规模如下：

\begin{table}[H]
\centering
\caption{数据概览}
\begin{tabular}{lr}
\toprule
指标 & 数值 \\
\midrule
总文件数 & 1521 \\
成功读取文件数 & 1521 \\
读取失败文件数 & 0 \\
空文件数 & 0 \\
总行数 & 3,040,479 \\
总列数 & 35 \\
\bottomrule
\end{tabular}
\end{table}

从基础读取结果看，所有已存在文件均能正常读取，没有发现空文件和读取失败文件，说明原始文件可读性较好。

\section{字段说明}

实际数据中的关键字段包括：

\begin{table}[H]
\centering
\caption{关键字段说明}
\begin{tabular}{ll}
\toprule
字段类型 & 字段 \\
\midrule
标识字段 & \texttt{date}, \texttt{time}, \texttt{sym} \\
价格字段 & \texttt{n\_close}, \texttt{n\_midprice}, \texttt{n\_bid1--n\_bid5}, \texttt{n\_ask1--n\_ask5} \\
数量字段 & \texttt{n\_bsize1--n\_bsize5}, \texttt{n\_asize1--n\_asize5} \\
成交字段 & \texttt{amount\_delta} \\
标签字段 & \texttt{label\_5}, \texttt{label\_10}, \texttt{label\_20}, \texttt{label\_40}, \texttt{label\_60} \\
EDA 辅助字段 & \texttt{source\_file}, \texttt{sym\_id}, \texttt{date\_id}, \texttt{session} 等 \\
\bottomrule
\end{tabular}
\end{table}

需要注意，实际数据中使用的是 \texttt{n\_close}，不是 \texttt{close}。因此后续代码和报告均应以实际字段名 \texttt{n\_close} 为准。

\section{文件完整性检查}

理论上，数据包含：

\begin{center}
10 只股票 $\times$ 79 个交易日 $\times$ 2 个交易时段 $=$ 1580 个文件组合
\end{center}

实际文件完整性检查结果如下：

\begin{table}[H]
\centering
\caption{文件完整性检查}
\begin{tabular}{lr}
\toprule
指标 & 数值 \\
\midrule
理论文件组合数 & 1580 \\
实际有效组合数 & 1521 \\
实际文件数 & 1521 \\
额外异常组合数 & 0 \\
缺失组合数 & 59 \\
\bottomrule
\end{tabular}
\end{table}

结果说明：理论上应有 1580 个股票-日期-session 组合，实际有效组合为 1521 个，缺失 59 个组合，未发现额外异常组合。缺失详情已保存至 \texttt{reports/missing\_files.csv}。

\section{基础数据质量检查}

\begin{table}[H]
\centering
\caption{基础数据质量检查}
\begin{tabular}{lr}
\toprule
检查项 & 结果 \\
\midrule
缺失值 & 未发现明显缺失 \\
重复行数 & 0 \\
Inf / -Inf 值 & 未发现 \\
读取失败文件 & 0 \\
空文件 & 0 \\
\bottomrule
\end{tabular}
\end{table}

整体来看，基础数据质量较好，可以进入进一步的订单簿结构检查和标签分析。

\section{订单簿异常检查}

\subsection{Bid-Ask 交叉异常}

\begin{table}[H]
\centering
\caption{Bid-Ask 交叉异常}
\begin{tabular}{lrr}
\toprule
异常项 & 数量 & 比例 \\
\midrule
\texttt{n\_bid1 > n\_ask1} & 30,853 & 1.01\% \\
\texttt{n\_bid2 > n\_ask2} & 32,468 & 1.07\% \\
\texttt{n\_bid3 > n\_ask3} & 33,806 & 1.11\% \\
\texttt{n\_bid4 > n\_ask4} & 34,696 & 1.14\% \\
\texttt{n\_bid5 > n\_ask5} & 35,674 & 1.17\% \\
\bottomrule
\end{tabular}
\end{table}

其中最重要的是 \texttt{n\_bid1 > n\_ask1}，约占全部样本的 1.01\%。正常订单簿中通常应满足 \texttt{bid1 <= ask1}，因此这部分样本可视为 crossed book 或盘口异常记录。

\subsection{Bid 档位内部顺序异常}

\begin{table}[H]
\centering
\caption{Bid 档位内部顺序异常}
\begin{tabular}{lrr}
\toprule
异常项 & 数量 & 比例 \\
\midrule
\texttt{n\_bid2 > n\_bid1} & 308 & 0.01\% \\
\texttt{n\_bid3 > n\_bid2} & 170 & 0.01\% \\
\texttt{n\_bid4 > n\_bid3} & 88 & 约 0.00\% \\
\texttt{n\_bid5 > n\_bid4} & 54 & 约 0.00\% \\
\bottomrule
\end{tabular}
\end{table}

\subsection{Ask 档位内部顺序异常}

\begin{table}[H]
\centering
\caption{Ask 档位内部顺序异常}
\begin{tabular}{lrr}
\toprule
异常项 & 数量 & 比例 \\
\midrule
\texttt{n\_ask2 < n\_ask1} & 1,307 & 0.04\% \\
\texttt{n\_ask3 < n\_ask2} & 1,168 & 0.04\% \\
\texttt{n\_ask4 < n\_ask3} & 802 & 0.03\% \\
\texttt{n\_ask5 < n\_ask4} & 924 & 0.03\% \\
\bottomrule
\end{tabular}
\end{table}

Ask 档位内部错序比例较低。所有 size 字段均未发现负数，说明挂单量字段没有明显非法负值。

\subsection{订单簿异常处理建议}

对于 \texttt{n\_bid1 > n\_ask1} 的样本，不建议在第一版建模前直接全部删除。更稳妥的策略是构造异常标记特征：

\begin{center}
\texttt{is\_crossed\_book = 1 if n\_bid1 > n\_ask1 else 0}
\end{center}

后续可以比较“保留异常标记”和“过滤异常样本”两种策略的验证集表现。

\section{Midprice 规则检查}

根据数据说明，若买一价和卖一价均不为 0，则理论中间价应为：

\begin{center}
\texttt{calc\_midprice = (n\_bid1 + n\_ask1) / 2}
\end{center}

如果一方为 0，则取非零价格；如果两方都为 0，则置为缺失。将该计算结果与数据自带的 \texttt{n\_midprice} 对比后得到：

\begin{table}[H]
\centering
\caption{Midprice 规则检查}
\begin{tabular}{lr}
\toprule
指标 & 数值 \\
\midrule
\texttt{midprice\_mismatch\_count} & 166,530 \\
\texttt{midprice\_mismatch\_ratio} & 5.48\% \\
\texttt{max\_abs\_midprice\_diff} & 0.0054545455 \\
\texttt{mean\_abs\_midprice\_diff} & 0.0000727803 \\
\bottomrule
\end{tabular}
\end{table}

结果说明：约 5.48\% 的样本中，数据自带 \texttt{n\_midprice} 与按 \texttt{n\_bid1}、\texttt{n\_ask1} 重新计算的中间价存在差异。但平均绝对差异较小，说明大多数不一致可能是轻微差异、数据源计算规则差异、浮点精度或特殊盘口状态导致。

后续建模建议优先使用数据自带的 \texttt{n\_midprice}，不要直接用重新计算值覆盖原字段。可以额外构造：

\begin{center}
\texttt{midprice\_diff = n\_midprice - calc\_midprice}
\end{center}

作为数据质量或盘口异常相关特征。

\section{标签方向检查}

由于项目说明中曾出现“当前 tick 相对于 N tick 之前”和“未来 N tick 相对于当前 tick”两种容易混淆的表述，因此本次 EDA 对标签方向进行了验证。

对每个文件单独检查，不跨文件计算，分别比较：

\begin{center}
未来方向：\texttt{n\_midprice[t+N] - n\_midprice[t]} \\
过去方向：\texttt{n\_midprice[t] - n\_midprice[t-N]}
\end{center}

检查结果如下：

\begin{table}[H]
\centering
\caption{标签方向检查}
\begin{tabular}{lrrl}
\toprule
标签 & Future 匹配率 & Past 匹配率 & 更符合方向 \\
\midrule
\texttt{label\_5} & 1.0000 & 0.6301 & future \\
\texttt{label\_10} & 1.0000 & 0.5554 & future \\
\texttt{label\_20} & 1.0000 & 0.5837 & future \\
\texttt{label\_40} & 1.0000 & 0.4504 & future \\
\texttt{label\_60} & 1.0000 & 0.3948 & future \\
\bottomrule
\end{tabular}
\end{table}

结论：五个标签均完全符合未来方向定义。因此，后续建模目标可以明确为：使用当前及过去若干 tick 的订单簿快照特征，预测未来 N tick 后中间价的移动方向。

\section{标签分布分析}

\begin{table}[H]
\centering
\caption{标签分布}
\begin{tabular}{lrrr}
\toprule
标签 & 下跌 0 & 不变 1 & 上涨 2 \\
\midrule
\texttt{label\_5} & 460,804 / 15.16\% & 2,122,348 / 69.80\% & 457,327 / 15.04\% \\
\texttt{label\_10} & 647,453 / 21.29\% & 1,762,358 / 57.96\% & 630,668 / 20.74\% \\
\texttt{label\_20} & 520,965 / 17.13\% & 2,006,167 / 65.98\% & 513,347 / 16.88\% \\
\texttt{label\_40} & 741,438 / 24.39\% & 1,586,710 / 52.19\% & 712,331 / 23.43\% \\
\texttt{label\_60} & 857,580 / 28.21\% & 1,361,229 / 44.77\% & 821,670 / 27.02\% \\
\bottomrule
\end{tabular}
\end{table}

\subsection{标签分布结论}

\begin{enumerate}
    \item \texttt{label\_5} 中不变类占比最高，达到 69.80\%，类别不均衡最明显；
    \item \texttt{label\_10} 的不变类占比下降到 57.96\%，上涨和下跌样本增加；
    \item \texttt{label\_20} 的不变类占比回升到 65.98\%，主要原因可能是 N=20 时标签阈值从 0.0005 提高到 0.001；
    \item \texttt{label\_40} 和 \texttt{label\_60} 的类别分布更均衡，尤其 \texttt{label\_60} 中上涨、下跌、不变三类比例相对接近；
    \item 后续训练时不能只看 accuracy，应重点关注下跌类 0 和上涨类 2 的 precision、recall 和 F0.5。
\end{enumerate}

\section{可视化结果与插图位置}

以下为建议插入的 EDA 图表位置。若图片文件已经位于 \texttt{reports/figures/}，可以将占位框替换为相应的 \texttt{\textbackslash includegraphics} 命令。

\begin{figure}[H]
\centering
\fbox{\parbox[c][5cm][c]{0.85\textwidth}{\centering 插图位置：各预测窗口标签分布图\\
\texttt{label\_5\_distribution.png}, \texttt{label\_10\_distribution.png}, \texttt{label\_20\_distribution.png}, \texttt{label\_40\_distribution.png}, \texttt{label\_60\_distribution.png}}}
\caption{不同预测窗口的标签分布}
\end{figure}

\begin{figure}[H]
\centering
\fbox{\parbox[c][5cm][c]{0.85\textwidth}{\centering 插图位置：各股票样本量分布\\
\texttt{samples\_per\_sym.png}}}
\caption{各股票样本量分布}
\end{figure}

\begin{figure}[H]
\centering
\fbox{\parbox[c][5cm][c]{0.85\textwidth}{\centering 插图位置：各交易日样本量分布\\
\texttt{samples\_per\_date.png}}}
\caption{各交易日样本量分布}
\end{figure}

\begin{figure}[H]
\centering
\fbox{\parbox[c][5cm][c]{0.85\textwidth}{\centering 插图位置：上午/下午样本量分布\\
\texttt{samples\_per\_session.png}}}
\caption{上午与下午 session 样本量分布}
\end{figure}

\begin{figure}[H]
\centering
\fbox{\parbox[c][5cm][c]{0.85\textwidth}{\centering 插图位置：中间价分布\\
\texttt{n\_midprice\_hist.png}}}
\caption{\texttt{n\_midprice} 分布}
\end{figure}

\begin{figure}[H]
\centering
\fbox{\parbox[c][5cm][c]{0.85\textwidth}{\centering 插图位置：成交金额变化分布\\
\texttt{amount\_delta\_hist.png}, \texttt{log\_amount\_delta\_hist.png}}}
\caption{\texttt{amount\_delta} 原始分布与对数变换分布}
\end{figure}

\begin{figure}[H]
\centering
\fbox{\parbox[c][5cm][c]{0.85\textwidth}{\centering 插图位置：Bid-Ask spread 分布\\
\texttt{spread\_hist.png}, \texttt{spread\_clipped\_hist.png}}}
\caption{Bid-Ask spread 原始分布与截尾分布}
\end{figure}

\begin{figure}[H]
\centering
\fbox{\parbox[c][5cm][c]{0.85\textwidth}{\centering 插图位置：相关性热力图\\
\texttt{correlation\_heatmap.png}}}
\caption{主要数值字段相关性热力图}
\end{figure}

\section{主要结论}

本次 EDA 得到以下结论：

\begin{enumerate}
    \item 数据整体可用。共读取 1521 个 CSV 文件，全部读取成功，无空文件、无读取失败文件。合并后共有 3,040,479 条样本和 35 个字段。
    \item 基础数据质量较好。未发现明显缺失值、重复行或无穷值。
    \item 文件组合存在缺失。理论上应有 1580 个股票-日期-session 组合，实际有效组合为 1521 个，缺失 59 个组合。后续划分训练集和验证集时应注意这些缺失交易组合。
    \item 订单簿整体结构基本合理，但存在少量异常。约 1.01\% 的样本存在 \texttt{n\_bid1 > n\_ask1}，即 crossed book 现象；bid/ask 档位内部错序比例较低；size 字段无负数。
    \item \texttt{n\_midprice} 与按 \texttt{n\_bid1}、\texttt{n\_ask1} 规则计算的中间价存在约 5.48\% 的不一致，但平均绝对差异较小。后续建议优先使用原始 \texttt{n\_midprice}，而不是直接覆盖。
    \item 标签方向已经确认。\texttt{label\_5}、\texttt{label\_10}、\texttt{label\_20}、\texttt{label\_40}、\texttt{label\_60} 均完全符合未来方向，即未来 N tick 的中间价相对于当前 tick 的移动方向。
    \item 标签存在类别不均衡。尤其 \texttt{label\_5} 和 \texttt{label\_20} 中不变类占比较高，后续模型不能只依赖 accuracy，应重点关注上涨和下跌类别的表现。
\end{enumerate}

\section{后续建议}

进入特征工程和建模阶段前，建议采取以下策略：

\begin{enumerate}
    \item \textbf{训练/验证划分：} 按 \texttt{date\_id} 做时间序列划分，不建议随机划分样本，避免未来信息泄露。
    \item \textbf{midprice 处理：} 优先使用原始 \texttt{n\_midprice}。如果需要使用重算中间价，应先确认其与标签生成逻辑的一致性。
    \item \textbf{异常订单簿处理：} 不建议第一版直接删除 crossed book 样本。可先构造 \texttt{is\_crossed\_book}、\texttt{spread}、\texttt{abs\_spread} 等特征，再比较过滤和不过滤的模型效果。
    \item \textbf{amount\_delta 处理：} 由于 \texttt{amount\_delta} 存在明显右偏和极端值，建议使用 \texttt{log1p(amount\_delta)} 或分位数截尾处理。
    \item \textbf{类别不均衡处理：} 训练分类模型时可使用类别权重、样本权重，或对上涨/下跌样本进行适度重采样。
    \item \textbf{评价指标：} 不应只看 accuracy。建议重点使用 precision、recall、F0.5，尤其关注标签 0 和标签 2 的预测效果。
    \item \textbf{特征工程方向：} 后续可构造 spread、order imbalance、bid/ask depth、rolling return、rolling volatility、过去 100 tick 的统计特征等。
\end{enumerate}

\end{document}
